\section*{Chapitre \ref{ch:m3:a-map-of-all-markets} --- Une carte de tous les marchés}

\begin{solution}{exo:m3:a-map-of-all-markets:1}
$1\,203.6/31\,800 = 3.8\%$ du stock en circulation chaque jour.
\end{solution}

\begin{solution}{exo:m3:a-map-of-all-markets:2}
New York est en retard de quatre heures sur l'UTC en été~: la séance principale court de 13:30--20:00 UTC et la fenêtre des taux de référence de 19:00--20:00 UTC.
\end{solution}

\begin{solution}{exo:m3:a-map-of-all-markets:3}
Par exemple~: les contrats à terme E-mini au pôle carnet d'ordres (carnet d'ordres centralisé, contrepartie centrale)~; les obligations d'entreprises américaines vers le pôle dealer
(demande de cotation (RFQ), règlement-livraison bilatéral, plateformes all-to-all en essor)~; un \omterm{def:m3:automated-market-makers:amm}{teneur de marché automatisé} (AMM) comme troisième type, un protocole
sans dealer ni chambre de compensation.
\end{solution}

\begin{solution}{exo:m3:a-map-of-all-markets:4}
$\lceil 24/7 \rceil = 4$ postes de huit heures~; $\lceil 24/11 \rceil = 3$ de douze.
\end{solution}

\begin{solution}{exo:m3:a-map-of-all-markets:5}
Les mesures diffèrent~: la BRI compte tous les instruments de change, nets des doubles comptes, sur un mois~; le chiffre des crypto-actifs porte sur une
journée, sur une seule plateforme et un seul instrument, déclaré par la plateforme elle-même~; et le volume des dérivés est exprimé en notionnel, non en valeur échangée.
\end{solution}

\begin{solution}{exo:m3:a-map-of-all-markets:6}
Ses positions sont tenues brièvement~: l'essentiel du volume provient des teneurs de marché et des traders à court horizon qui font tourner
leur inventaire, et non de détenteurs d'opinions pendant des semaines.
\end{solution}

\begin{solution}{exo:m3:a-map-of-all-markets:7}
Jour ouvré~: postes 00--09, 08--17 et 16--01 UTC avec 2, 3 et 3 traders, soit 8 postes de trader~; dimanche~: les mêmes postes avec un trader chacun, soit 3.
\end{solution}

\begin{solution}{exo:m3:a-map-of-all-markets:8}
La taille n'est pas synonyme de bénéfice~: les plus grands marchés sont les plus concurrentiels et servis par les plus gros bilans~; l'opportunité se trouve là où les compétences d'une société
rencontrent des participants qui négocient pour d'autres raisons que le prix, à une échelle qu'elle peut servir.
\end{solution}

\begin{solution}{pb:m3:a-map-of-all-markets:1}
\textbf{1.} Trois postes de neuf heures se chevauchant d'une heure. \textbf{2.} 00:00--09:00, 08:00--17:00 et 16:00--01:00 UTC. \textbf{3.} 2, 3 et 3
(crypto-actifs et électricité la nuit à partir de 05:00~; crypto-actifs, électricité et États-Unis le jour~; crypto-actifs et États-Unis le soir). \textbf{4.} Un par poste. \textbf{5.} $5
\times 8 + 2 \times 3 = 46$. \textbf{6.} Premier poste~: financement à 00:00 et à 08:00 avec l'échéance des options (dans le chevauchement)~; deuxième~: à 08:00,
l'enchère day-ahead à 10:00, les \omterm{def:m3:metals:official}{prix officiels du LME} et l'ouverture américaine, financement à 16:00 (dans le chevauchement)~; troisième~: 16:00, la fenêtre des taux de référence
et la clôture américaine à 20:00, et financement à minuit (dans le chevauchement). \textbf{7.} Le deuxième. \textbf{8.} Le troisième. \textbf{9.} Positions ouvertes et
ordres, limites utilisées, virements en attente et appels de marge, incidents, et les moments fixes du poste suivant. \textbf{10.} En hiver, la session américaine
et la fenêtre des taux de référence sont décalées d'une heure en UTC, de même que l'enchère européenne, si bien que le contenu des deuxième et troisième postes change, et
il faut exprimer à nouveau les fenêtres de desk en heure locale. \textbf{11.} Deux postes laisseraient deux heures non couvertes avec un chevauchement d'une heure
($\lceil 24/11 \rceil = 3$)~: des postes de douze heures en demandent encore trois, au prix de la fatigue. \textbf{12.} Oui, pour couvrir les crypto-actifs la nuit et
l'électricité tôt le matin, au prix d'un second bureau et de passations entre fuseaux horaires. \textbf{13.} Automatisation des cotations dans des limites, couverture et surveillance~;
les humains doivent surveiller les limites de risque, les incidents, les problèmes de plateforme et les mouvements importants. \textbf{14.} Un planning crypto réduit, avec automatisation et un humain d'astreinte.
\textbf{15.} 46 postes de trader par semaine, soit environ neuf traders et demi à temps plein à cinq postes chacun, multipliés par le coût indiqué, auxquels s'ajoutent le chevauchement et
la couverture des congés. \textbf{16.} Celui dont le profit par heure de trader est le plus faible, mesuré et non supposé~; la crypto coûte le plus cher en personnel, le bureau américain le moins.
\textbf{17.} Vers des marchés qui réutilisent la même technologie et les mêmes horaires (par exemple les dérivés de crypto-actifs sur des plateformes régulées) ou comblent les heures creuses.
\textbf{18.} Les marchés à forte \omterm{def:m3:a-map-of-all-markets:velocity}{vitesse de rotation} récompensent la rapidité de la société et la rotation de son inventaire~; ceux à faible vitesse exigent un bilan qu'elle n'a peut-être pas.
\textbf{19.} \emph{Résultat nommé~:} trois postes de neuf heures par jour (00--09, 08--17, 16--01 UTC) avec 2, 3 et 3 traders les jours ouvrés et un seul le week-end.
\textbf{20.} Pour choisir où négocier, et voir ce que chaque marché exige avant d'y aller.
\end{solution}

\begin{solution}{iq:m3:a-map-of-all-markets:1}
Le change (FX) d'abord, puis les bons du Trésor américain, puis les actions et les crypto-actifs~; le classement masque des mesures différentes (notionnel ou valeur, net ou brut, une plateforme ou toutes),
la \omterm{def:m3:a-map-of-all-markets:velocity}{vitesse de rotation}, la part de ce volume à laquelle une société peut accéder, et les marges réalisées.
\iqlookfor{mesures, \omterm{def:m3:a-map-of-all-markets:velocity}{vitesse de rotation} et accessibilité.}
\end{solution}

\begin{solution}{iq:m3:a-map-of-all-markets:2}
Positions ou ordres que le desk suivant ne comprend pas, limites incohérentes d'un fuseau horaire à l'autre, événements survenus dans l'intervalle entre postes, différences de pratiques
locales, fatigue~; des listes de contrôle lors de la passation, un seul registre faisant foi et le chevauchement les réduisent.
\iqlookfor{perte d'information lors de la passation.}
\end{solution}

\begin{solution}{iq:m3:a-map-of-all-markets:3}
Volume que la société pourrait négocier, multiplié par la fourchette qu'elle pourrait capter nette de sélection adverse et de frais, à partir de la profondeur du marché, de ses participants et de la concurrence,
vérifié par rapport aux profits des sociétés déjà présentes, et minoré du coût d'accès.
\iqlookfor{volume multiplié par la captation nette, et concurrence.}
\end{solution}

\begin{solution}{iq:m3:a-map-of-all-markets:4}
Un marché où les participants négocient pour d'autres raisons que le prix, où la technologie compte plus que le bilan et où l'accès est peu coûteux~: par exemple la tenue de
marché dans un jeune marché de dérivés, avec les réserves liées au risque de plateforme.
\iqlookfor{une adéquation raisonnée entre l'edge et les contraintes.}
\end{solution}

\begin{solution}{iq:m3:a-map-of-all-markets:5}
Un service de calendrier qui tient à jour, pour chaque marché, les règles de séance en heure locale, les jours fériés et les fermetures anticipées, qui convertit en UTC, pour toute date, à partir des données de fuseaux horaires, et expose
les requêtes «~ouvert maintenant~» et «~prochain événement~», versionné et testé par rapport aux calendriers publiés des marchés.
\iqlookfor{règles locales, conversion de fuseaux horaires et versionnement.}
\end{solution}

\begin{solution}{iq:m3:a-map-of-all-markets:6}
Choisir des marchés qui partagent une infrastructure (tenue de marché en carnet d'ordres), couvrent des horaires complémentaires et présentent une faible corrélation des profits~: par exemple les contrats à terme sur indices
actions, les \omterm{def:m3:perpetual-futures:perp}{contrats perpétuels} sur crypto-actifs et l'intrajournalier européen de l'électricité~; justifier par l'edge, les horaires et le capital.
\iqlookfor{technologie partagée, horaires complémentaires, profits diversifiés.}
\end{solution}
